\section{绪论}

\subsection{研究背景}

近十年来，视觉语言模型取得了长足进步。以 CLIP 为代表的对比式模型通过数亿图文对的联合训练，首次证明视觉和语言可以在统一嵌入空间中实现有效对齐。在此基础上，BLIP-2、LLaVA、Qwen-VL 等生成式视觉语言模型进一步将视觉编码器与大语言模型连接，实现了图像描述、视觉问答和跨模态推理等功能。这些模型已被广泛应用于自动驾驶、医疗影像辅助分析、智能监控和增强现实等实际领域。

然而，上述模型的性能高度依赖输入图像的质量。在现实环境中，图像常因光照、天气、传感器老化、传输压缩或遮挡等原因出现视觉退化。当图像分辨率下降、噪声增大或关键区域被遮挡时，视觉编码器提取的语义特征会变得不可靠，下游的类别判断和文本生成可能产生明显偏差。例如在雨雾天气下，自动驾驶系统对交通标志的识别能力会显著降低。提升退化条件下的多模态理解鲁棒性是视觉语言模型走向实际应用面临的关键问题之一。

另一方面，脑电图作为一种无创的脑机接口技术，能够以毫秒级时间分辨率记录大脑皮层在视觉刺激下的电生理活动。近年来，EEG 视觉语义解码研究取得了一定进展：已有工作表明，从被试观看自然图像时的 EEG 信号中，可以解码物体类别、检索语义相关图像，甚至与预训练视觉语言模型的嵌入空间建立映射关系。这引入了一个值得研究的问题：当视觉输入质量受损时，观察者大脑中形成的神经表征是否能为机器视觉系统提供额外的语义线索。

本课程设计正是在上述交叉背景下展开，研究视觉退化条件下真实配对 EEG 信号对视觉语言模型的语义辅助作用。需要明确的是，本设计的研究定位并非``从脑电直接读取语言''，而是将 EEG 作为一种视觉语义增强信号，在所有实验结论中，均通过与 scrambled EEG 和随机 EEG 的对比来给出客观判断依据。

\subsection{课程设计目标与意义}

本课程设计的目标是：在 Thought2Text 小样本 EEG-image 配对数据集上，构建视觉退化条件下的 EEG 辅助视觉语言理解系统，系统评估真实 EEG 能否在六种退化场景下改善语义识别和图像描述生成的表现。

课程设计的实践意义涵盖三个层面：

\textbf{方法实践层面}：课题综合运用了深度学习模型设计（EEG encoder、gated fusion、cross-attention）、预训练模型的使用与适配（CLIP 视觉和文本编码器、Qwen2-VL 生成模型）、参数高效微调（LoRA）、对比学习（InfoNCE）、图像增强与退化建模等相关技术。在约两个月的项目周期内，完成了从文献调研、数据组织、模型设计、代码实现、训练调参到对照评估的全流程实践。

\textbf{系统构建层面}：本设计从受约束的语义分类任务逐步过渡到生成式描述任务，先通过 A2 semantic fusion 在 CLIP 空间中验证 EEG 语义辅助的有效性，再通过 VTF token-level fusion 将 EEG 信息推进到视觉 token 层面，最终接入 Qwen2-VL 生成式模型完成自由 caption 输出。这一递进式构建思路有助于在每一步都有可验证的中间结果。

\textbf{科学验证层面}：本设计强调严格的对照实验——real EEG、shuffled EEG、random EEG 和 vision-only 四组对比贯穿所有评估，确保任何性能提升都能归结于真实配对 EEG 携带的语义信息，而非模型容量增加带来的噪声拟合。

\subsection{研究内容与技术路线}

本课程设计主要完成以下四项工作：

\begin{enumerate}
  \item 构建视觉退化条件下 EEG 辅助语义识别与生成式描述的完整任务框架，包含六种退化条件（clean、lowres16、strong\_blur、strong\_noise、occlusion50、mixed）和四种对照模式，采用 image-level split 防止训练-测试信号泄漏；
  \item 设计并实现 A2 temporal-spectral-spatial EEG 语义融合模型，在 CLIP 语义空间中同时利用 EEG 的时间动态、频谱分布和空间通道关系，通过门控残差融合对退化图像语义表征进行补充；
  \item 提出 VTF token-level 融合方法，将 EEG token 通过 visual-to-EEG cross-attention 注入 CLIP ViT visual token 序列，生成增强视觉表示，为后续生成式模型提供 EEG-conditioned 视觉输入；
  \item 构建基于 Qwen2-VL 的生成式视觉语言模型，将增强视觉 token 输入预训练生成模型，使用 LoRA 进行参数高效微调，在六种退化条件下产生自由形式图像描述，并通过 best-of-N reranking 提升输出质量。
\end{enumerate}

技术路线的总体思路是：先利用预训练 CLIP 的语义空间作为跨模态桥梁，降低对 EEG 数据量的依赖；再在分类任务上严格验证 EEG 语义增强的效果；然后逐步推进到 token-level fusion 和生成式输出。这一由简到繁、从定量到生成的分层设计，使得每一步都有可控的评估和明确的验证依据。

\subsection{报告结构安排}

本报告共分为十个章节。第 1 章为绪论；第 2 章介绍相关技术与理论基础；第 3 章介绍数据集与任务定义；第 4 章阐述系统总体设计；第 5 章详细描述 A2 EEG 语义融合模型；第 6 章介绍 token-level EEG-enhanced VLM 融合方法；第 7 章阐述生成式视觉语言模型的设计与实现；第 8 章给出实验结果与分析；第 9 章总结工程实现与代码结构；第 10 章进行总结与展望。
